\section{The VCVio Library}\label{sec:vcvio}

This library provides a formal framework for reasoning about computations that make \emph{oracle
queries}. Many cryptographic primitives and interactive protocols use oracles to model (or simulate)
external functionality such as random responses, coin flips, or more structured queries. The VCVio
library "lifts" these ideas into a setting where both the abstract specification and concrete
simulation of oracles may be studied, and their probabilistic behavior analyzed.

The main ingredients of the library are as follows:

\begin{definition}[Specification of Oracles]
    \label{def:oracle_spec}
    An oracle specification on an index type $\iota$ is a map $\mathsf{spec} : \iota \to \mathsf{Type}$.
    Each $t : \iota$ is a possible query, and $\mathsf{spec}\ t$ is the type of answers to that
    query. The query type is therefore the index type itself (\verb|spec.Domain| is $\iota$), and the
    answer type depends on the query (\verb|spec.Range t| is $\mathsf{spec}\ t$). An oracle
    specification is the same data as a polynomial functor with positions $\iota$ and directions
    $\mathsf{spec}\ t$. The specification only describes interfaces, not how the oracles behave.
    \lean{OracleSpec}
    \leanok
\end{definition}

\begin{definition}[Standard Oracle Specifications]
    \label{def:standard_oracle_specs}
    \begin{itemize}
        \item $A \to_{\mathsf{o}} B$: oracles that take inputs in $A$ and answer every query in $B$.
        \item \verb|emptySpec|, written $[]_{\mathsf{o}}$: no oracles.
        \item \verb|coinSpec| $= \mathsf{Unit} \to_{\mathsf{o}} \mathsf{Bool}$: a coin-flipping oracle.
        \item \verb|unifSpec|: indexed by $n \in \mathbb{N}$, the query $n$ is answered in
        $\mathsf{Fin}\ (n+1) = \{0, \ldots, n\}$; with uniform answers it samples uniformly from this set.
    \end{itemize}
    Specifications combine by disjoint sum $\mathsf{spec}_1 + \mathsf{spec}_2$.
    \lean{OracleSpec.emptySpec, coinSpec, unifSpec}
    \leanok
    \uses{def:oracle_spec}
\end{definition}

We often require extra properties of query and answer types, such as decidable equality or
finiteness. These are stated as ordinary instance hypotheses on the index type and on
\verb|spec.Range t| (quantified over $t$), rather than as bundled classes on the specification.

\begin{definition}[Oracle Computation]
    \label{def:oracle_computation}
    An oracle computation \verb|OracleComp spec| $\alpha$ is a program that makes queries to the oracles of
    \verb|spec| and returns a value of type $\alpha$. It is the free monad of the polynomial functor
    of \verb|spec|, so every computation either
    \begin{itemize}
        \item returns a value without making any queries (\verb|pure|), or
        \item makes one query $t$ and continues with a function of the answer in \verb|spec.Range t|
        (\verb|queryBind|).
    \end{itemize}
    There is no failure constructor. Computations that may fail are written in
    \verb|OptionT (OracleComp spec)|. Computations whose only oracle is \verb|unifSpec| are called
    probabilistic computations, \verb|ProbComp|.
    \lean{OracleComp}
    \leanok
    \uses{def:oracle_spec}
\end{definition}

\begin{definition}[Handling Oracle Queries]
    \label{def:handling_oracle_queries}
    To run oracle computations, we need a way to handle (or implement) the oracle queries.
    A query implementation \verb|QueryImpl spec m| maps every query $t$ to a computation in a monad
    $m$ returning an answer in \verb|spec.Range t|. Depending on the monad, this allows various
    interpretations of the queries.
    \lean{QueryImpl}
    \leanok
    \uses{def:oracle_spec}
\end{definition}

\begin{definition}[Probabilistic Semantics of Oracle Computations]
    \label{def:probabilistic_semantics_of_oracle_computations}
    A computation denotes a subprobability measure on its outputs, \verb|evalDist mx|. The missing
    mass represents failure or nontermination; the output distribution does not add an
    \verb|Option| outcome. For oracle computations, the semantics is fixed by choosing a probability
    measure on the answers to each query. The usual choice is the uniform measure, as for
    \verb|ProbComp|. Probabilities of outputs and events are defined from this measure.
    \lean{evalDist}
    \leanok
    \uses{def:oracle_computation}
\end{definition}

\begin{definition}[Simulating Oracle Queries with Other Oracles]
    \label{def:sim_oracle}
    Given a query implementation of \verb|spec| in a monad $m$, \verb|simulateQ| runs an oracle
    computation by substituting the implementation for each query. Taking $m$ to be a computation
    with access to other oracles, possibly with state (for example a state monad over \verb|OracleComp spec'|),
    simulates one specification by another.
    \lean{simulateQ}
    \leanok
    \uses{def:oracle_computation, def:handling_oracle_queries}
\end{definition}

\begin{definition}[Logging \& Caching Oracle Queries]
    \label{def:logging_caching_oracle_queries}
    Using the simulation framework, we can add logging and caching behaviors to oracle queries:
    \begin{itemize}
        \item Logging records all queries made during a computation
        \item Caching remembers query responses and reuses them for repeated queries
    \end{itemize}
    Logging forwards every query and records it. Caching forwards only cache misses and answers
    repeated queries from the cache.
    \lean{OracleSpec.loggingOracle, OracleSpec.cachingOracle}
    \leanok
    \uses{def:oracle_computation, def:sim_oracle}
\end{definition}

\begin{definition}[Random Oracle]
    \label{def:random_oracle}
    A random oracle is implemented as a caching oracle that uses lazy sampling:
    \begin{itemize}
        \item On first query: generates a uniform random response and caches it
        \item On repeated queries: returns the cached response
    \end{itemize}
    \lean{OracleSpec.randomOracle}
    \leanok
    \uses{def:oracle_computation, def:handling_oracle_queries}
\end{definition}
